\label{sec:human_eval}
To complement our automatic evaluations and better understand the effectiveness and limitations of \model, we conducted human evaluations. This study involved over 100 literature review questions and more than 15 participants, including Ph.D. students, research scientists, and university professors with expertise in the relevant fields. In total, we curated more than 400 fine-grained expert evaluations on human and model answers. 


\subsection{Human Evaluation Design}
\paragraph{Evaluations against human experts.}
For human evaluations, we use 108 question-answer pairs from \textsc{ScholarQA-Multi}, written by experts. 
We run three models on these questions to generate answers with citations: GPT4o (without external retrieval), \model with GPT4o as the generator (OS-GPT4o), and \model with our trained 8B model (OS-8B). Expert annotators are then asked to evaluate the model-generated answers against human-written answers.

Each evaluation involves presenting a question, a model-generated answer, and a human-written answer. Expert annotators then conduct fine-grained assessments of each answer and provide pairwise preference judgments between the two. 
For fine-grained evaluations, we use the five-scale evaluation criteria described in \Sref{sec:dataset} (\covh, \orgh, \relh), with annotators scoring both model and human answers using the same rubrics. 
For usefulness (\useh), annotators assign scores on a scale from 1-5, which we convert into three classes: Not Useful (1-2), Neutral (3), and Useful (4-5). We then calculate the percentage of answers that fall into the Useful category. For pairwise preference, annotators either choose one of the answers or mark a ``tie'' if they judge both answers to be of equal quality. Optionally, experts provide explanations on why one answer is better than the other.  

\if0{
We compare the following models: 
\begin{itemize}[leftmargin=10pt]
    \item {\bf \model with GPT4o}: Answers generated by our \model using GPT4o as the generator LM. 
    \item {\bf \model with Our 8B}: Answers generated by our \model pipeline using our newly trained 8B model, which is fully open. 
\end{itemize}
}\fi

\paragraph{Expert annotators for answer writing.} 
Our expert annotators for question and answer writing are 12 Ph.D. students and post-doctoral researchers from research institutions across the United States, all of whom have at least three years of research experience and published multiple papers in journals or conferences from their fields.  
Our annotators' expert areas span computer science (Natural Language Processing, computer Vision, Human-Computer Interaction), physics (Astrophysics and Photonics / Optics), and biomedical (Neuroscience, Bioimaging) domains, and we assign our expert annotators to  questions in their expertise. On average, we paid 35-40 USD per person. 

\paragraph{Expert annotators for evaluations.} 
A total of 16 expert annotators from the three fields contributed to our evaluations, with 12 of them also participating in answer generation. All expert annotators met the same qualifications as those who composed the answers. To minimize potential biases, we ensured that annotators did not evaluate responses to their own questions by assigning evaluation tasks to different groups of experts. 
Each instance was reviewed by 1 to 3 expert annotators, depending on availability. The inter-annotator agreement was 0.68 using pairwise comparison with ties, and 0.70 using a relaxed approach, wherein ties were merged. On average, each expert spent five minutes per instance on evaluation, and received compensation ranging from 25–35 USD. 

\begin{table}[t!]
\small
    \centering
    \begin{tabular}{l |ccc|c|ccc}
    \toprule
     & \multicolumn{3}{c}{Fine-grained (1-5, Avg.)} & Overall Usefulness &\multicolumn{3}{c}{Relative to Human (\%)} \\
        &  \orgh & \covh & \relh & \useh~(\%)  & Win & Tie & Lose  \\
    \midrule
    GPT4o & 4.63 {\textcolor{osgreen}{\bf (+0.4)}} & 4.06 {\textcolor{orange}{\bf (-0.2)}} & 4.50 {\textcolor{orange}{\bf (-0.1)}} & 69.7 {\textcolor{orange}{\bf (-13.9)}} & 31.9 & 13.8 & \bf 54.2  \\
    \textsc{OS}-8B & 3.82  {\textcolor{orange}{\bf (-0.3)}} & 4.30 {\textcolor{osgreen}{\bf (+0.7)}} & 4.00 {\textcolor{orange}{\bf (-0.4)}} &72.1 {\textcolor{osgreen}{\bf (+8.7)}} & \bf 50.8 & 12.3 & 36.9 \\
    \textsc{OS}-GPT4o & {4.47} {\textcolor{osgreen}{\bf (+0.8)}} & { 4.38} {\textcolor{osgreen}{\bf (+0.9)}}  & { 4.30} {\textcolor{gray}{\bf (0.0)}}  & { 80.0} {\textcolor{osgreen}{\bf (+22.5)}}  & \bf 70.0 & 6.8 & 23.2 \\
    \bottomrule
    \end{tabular}
        \caption{{\bf Human evaluation results}. Fine-grained aspect evaluations are conducted on a five-point scale across four aspects using our detailed instructions and rubrics. Values in parentheses represent the relative performance difference; \textcolor{osgreen}{\bf (+)} indicates that the model shows higher performance, and \textcolor{orange}{\bf (-)} indicates that the human shows higher performance. 
    }
    \label{table:human_evaluation_results}
\end{table}
\subsection{Human Evaluation Results}
\begin{figure*}[t!]
    \centering
    \includegraphics[width=\textwidth]{figs/human_eval.pdf}
    \vspace{-4mm}
    \caption{{\bf Fine-grained evaluation results. } Score distributions between 1) GPT4o (top), \model with 8B (middle), \model with GPT4o with Human (bottom). 
    }
    \label{fig:human_eval}
\end{figure*}

\paragraph{Results of human evaluations.}
Table~\ref{table:human_evaluation_results} presents the average scores for each evaluation aspect, alongside the relative win rates against human responses. 
Figure~\ref{fig:human_eval} illustrates the score distributions for Human, GPT4o, and \model with Llama 3 8B and GPT4o. 
Notably, both OS-GPT4o and our OS-8B versions outperform human answers in over 50\% of cases, with their advantage primarily attributed to their ability to provide a greater breadth and depth of information (coverage; \covh). In contrast, GPT4o, which lacks retrieval capabilities, demonstrates significantly limited coverage and wins in fewer than 35\% of cases, with its overall usefulness rated much lower than responses from humans and the other two models. These results highlight that even for state-of-the-art models, synthesizing and answering scientific literature review questions remains a challenging task, consistent with our findings on \data. 
Overall, \model-GPT4o and \model-8B are rated as Useful in 80\% and 72\% of the queries, respectively. 

While \model with a smaller open 8B LM already suppresses human experts, the 8B model's output is judged to be less organized or fluent than the current state-of-the-art private LM-based \model. We found that GPT4o incorporates feedback more effectively, and generally generates longer and more fluent outputs, leading to significantly higher organization scores compared to 8B based \model as well as humans. 


\paragraph{Effects of length control on model responses. }
While we found model outputs to often be preferred over human outputs, one potential confounding factor is the significant difference in their output length--\model-GPTo and \model-8B are 2.4 times and 2.0 times longer than human-written answers, respectively, which may affect human judgment~\cite{dubois2024length}. 
To understand the effect of output length, we conducted a controlled experiment: for randomly sampled 50 questions, we generate shorter version of responses for \model-GPT4o, by prompting GPT4o to create summaries of the responses that fall under 300 words. As a result, we collected \model answers that average around 333 words, which is close to the average human answer length. 
We then conduct the same human evaluation and evaluate fine-grained and overall responses. 
On average, the shortened GPT4o scores 4.5 for organization, 4.6 for coverage, and 4.6 for relevance. 
The shortened \model-GPT4o responses are preferred or tied with expert answers in 75\% of the queries. 
The experimental results show that the model's superior performance is not merely due to the increased length of the \model answers. Moreover, human annotators' explanations often mention that both shortened \model and human answers could be improved by incorporating more details, implying that a 300-word restriction may limit the utility of answers.  



\paragraph{Analyses on human explanations for pair-wise explanations.}  
We randomly sampled 59 instances with free-form explanations of pairwise preferences and conducted a manual analysis to identify factors that influence overall preferences. Specifically, we examined whether the explanations referenced one or more of the following four categories: organization, relevance, coverage, and citations. While the first three categories align with the fine-grained human evaluation criteria, the citation category also considers the quality of the cited papers (e.g., whether the system includes key representative papers in the field). 
Our analysis revealed that 12\%, 23\%, 29\%, and 9\% of the explanations cited organization, relevance, coverage, and citations, respectively, as key factors in pairwise decisions. This suggests that coverage plays a crucial role in how humans assess the quality of responses, with annotators largely favoring model-generated answers for their greater coverage and depth of information. However, annotators also noted that the citations provided by models could be improved, pointing out that the suggested papers were occasionally outdated or less relevant compared to more representative work. Appendix~\ref{table:explanationst} shows example explanations. 
